\chapter{Abstract}

\textbf{Abstract:}

In this project, the conceptual and theoretical design, as well as the construction of a semiautomatic system for generating compost from organic waste through an In-Vessel composting process, were carried out. This process is characterized by the use of a closed container.

In this process, it is essential to consider the evolution of the environment inside the container in order to maintain favorable conditions to promote bacterial growth.

The bacterial culture is going to provide different chemical and physical properties and turn them into natural soil fertilizer.

This project has as its main aim to build up a mechatronic system that comes into being from the work carried out during the research methodology and Terminal Work I.

In this way, this report will provide an exhaustive overview of each phase that marked the conception and execution of this project. Over the course of this path, the challenges presented at each stage of implementation will be addressed in detail.

All solutions and strategies to raffle each of the obstacles will be examined; furthermore, the acquired knowledge and the evolution that the team has presented during the development of the project will be addressed in detail. This analysis will not just provide a clear view of the implementation of the idea, but it is also expected to serve as a valuable source of learning for future projects.


%\textbf{Palabras Clave:} Compostaje, Residuos alimentarios, Sistema semiautomático, Sistema mecatrónico, Proceso In-Vessel, Fertilizante orgánico, Crecimiento bacteriano, Monitoreo ambiental, Compost, Tratamiento de residuos orgánicos. \\

%\textbf{Keywords:} Composting, Food Waste, Semiautomatic System, Mechatronic System, In-Vessel Process, Organic Fertilizer, Bacterial Growth, Environmental Monitoring, Compost, Organic Waste Treatment. \\

\textbf{Keywords:} Compost, Composting, Organic Waste, Semiautomatic System, Mechatronic System, In-Vessel Process, Bacterial Growth, Closed Container. 

%%%%%%fin del archivo
\endinput 